\chapter{Column distributions}\label{ch:gutter}
\begin{goals}
\textbf{In this chapter} you learn to read the small drawing above each column header, and to use it:
hover to read a value, click to filter, drag to filter a range, double-click for the full statistics.
\end{goals}

Above every column header lamina draws the column's distribution: a histogram for a number column, a bar
of its most common values for a text one, and a thin bar showing how much of it is blank. They are
drawn when the file opens, so the first thing you see of a new file is the shape of every column in it.
\menu{View>Column distributions} turns them off and on.

\shot{gutter}{The distributions over the sample's columns: histograms for the grades (log-scaled, marked
\ui{log}), a bar of the four lithologies over \col{LITO}, and flat glyphs over the constant block sizes.}

\section{What they show}\label{sec:glyphs}

\begin{description}
\item[A histogram] for a number column: 22 bars between the column's minimum and maximum.
\item[A top-values bar] for a text column: the eight most common values as coloured segments, each as
wide as its share.
\item[The blank bar] under either: its share of the column that is empty.
\item[\ui{≈}] the glyph is drawn from a sample, not the whole file (below).
\item[\ui{log}] the histogram is drawn on a log scale of value.
\end{description}

\subsection{The sample}\label{sec:gutter-sample}

On a file of more than 8{,}192 rows, the distributions are drawn from a sample of 8{,}192 rows: 48 short
runs of consecutive rows, spread evenly from the start of the file to its end. Spread, because a file is
often ordered (a block model by elevation, an assay table by hole), and the first rows alone would show
one corner of it. The sample is read when the file opens, a few blocks of the file, and costs little at
any size.

A sample can miss rare values: a handful of high grades in ten million blocks may not be in it, and the
histogram's maximum is then the sample's. For exact figures, double-click the glyph (\cref{ch:stats}).

\subsection{Log scale}\label{sec:auto-log}

Grades are usually skewed: most values small, a long tail of large ones. On a linear histogram the tail
squeezes everything else into the first bar. So when a number column is skewed, lamina draws it on a
log scale and marks it \ui{log}.

The test is made on the sample's positive values (zeros are common in a grade column, and are set aside;
more than 2\% negative values rules a log scale out): when their skewness is above 1.5 and taking logs
reduces it by more than half a unit, the column is drawn on a log scale. \menu{View>Auto log-scale skewed
columns} turns the test off. For one column, the header's menu has \ui{Log distribution} to switch it by
hand.

\section{With a filter on}\label{sec:gutter-filtered}

When a filter is on, each glyph shows the matching rows' distribution over the whole column's, so you
see where the filtered rows sit in each column: filter to \code{LITO = "SULPHIDE"} and the \col{Z}
histogram shows the sulphide at depth, the \col{Cu\_pct} one its share of the high grades.

\section{Using them}\label{sec:gutter-use}

\begin{description}
\item[Hover] shows the value under the pointer: on a histogram, the value at that point of the axis
(\code{≈ 0.42}); on a top-values bar, the value and its share of the sample (\code{SULPHIDE · 27\%}).
\item[Click] filters to it: a histogram bar gives \code{Cu\_pct between 0.36 and 0.73}, a text
segment \code{LITO == "OXIDE"}. The filter box shows the expression, so you can edit it.
\item[Drag] across a histogram filters to the range you drag over; across a top-values bar, to the
values it covers (\code{LITO in ("SULPHIDE", "TRANSITION")}). While you drag, a tooltip shows the range.
\item[Double-click] opens the column's statistics (\cref{sec:one-column}).
\end{description}

On a log-scaled histogram, a click or a drag is mapped through the log scale, so the range you get is the
one you see. The numbers are rounded to a few significant figures, which is what you would have typed.

\subsection{When the filter runs}\label{sec:brush-modes}

A click or a drag on a large file starts a scan, and dragging a slightly different range would start
another. \menu{Data>Gutter brush} chooses what happens:

\begin{description}
\item[\ui{Auto (apply if small, stage if huge)}] the default. On a view of fewer than 2~million rows the
filter applies when you let go. On a larger one, the expression is written into the filter box and the
status line says \code{filter staged — press Enter to apply}: adjust it if you like, then press
\key{Enter}.
\item[\ui{Apply on release}] always apply when you let go.
\item[\ui{Stage (press Enter to run)}] never apply by itself.
\end{description}

The setting governs clicks and drags on histograms. A click or a drag on a top-values bar always applies
at once.

\begin{tryit}
Find where the sulphide sits.
\begin{enumerate}
\item Open the sample. Hover over the top-values bar above \col{LITO}: four segments, WASTE the widest.
\item Click the SULPHIDE segment. The filter box reads \code{LITO == "SULPHIDE"} and the status line
counts the matches.
\item Look at the \col{Z} histogram: the matching rows are drawn over the column's, all at depth.
\item Drag across the upper half of the \col{Cu\_pct} histogram. The filter becomes a range of copper.
It replaces the lithology filter: a brush writes a new filter rather than adding to the old one. To keep
both, edit the filter box: \code{LITO = "SULPHIDE" and Cu\_pct > 1}.
\end{enumerate}
\end{tryit}
